\documentclass[12pt]{article}
\usepackage{amsmath, amssymb, amsthm}
\usepackage[utf8]{inputenc}
\usepackage[T1]{fontenc}
\usepackage{geometry}
\usepackage{hyperref}

\geometry{a4paper, margin=2.5cm}

\newtheorem{lemma}{Lemma}
\newtheorem{theorem}{Theorem}
\theoremstyle{remark}
\newtheorem*{remark}{Remark}

\title{\textbf{Bounds on the Number of Divisors}}
\author{0xHaru}
\date{\today}

% ---------------------------------------------------------------
\begin{document}
\maketitle

% ===============================================================
\section*{Preliminary Lemma}
% ===============================================================

\begin{lemma}
Let $d \in \mathbb{N}$ be a divisor of $n \in \mathbb{N}$.  Then at least one of $d$
or $k = n/d$ is less than or equal to $\sqrt{n}$.
\end{lemma}

\begin{proof}[Proof 1 (by contradiction)]
Since $d$ divides $n$, there exists a natural number $k$ such that $d \cdot k = n$.

Suppose for contradiction that \emph{both} $d$ and $k$ are strictly greater
than $\sqrt{n}$:
\[
    d > \sqrt{n} \quad \text{and} \quad k > \sqrt{n}.
\]
Multiplying these two inequalities gives
\[
    d \cdot k > \sqrt{n} \cdot \sqrt{n} = n,
\]
which contradicts $d \cdot k = n$.  Therefore, at least one of $d \leq \sqrt{n}$
or $k \leq \sqrt{n}$ must hold.
\end{proof}

\smallskip

\begin{proof}[Proof 2 (direct)]
Since $d$ divides $n$, there exists a natural number $k = n/d$ such that
$d \cdot k = n$.  We consider two cases.

\medskip
\noindent\textit{Case $d \leq \sqrt{n}$.}\;
The thesis holds immediately, since $d$ itself is at most $\sqrt{n}$.

\medskip
\noindent\textit{Case $d > \sqrt{n}$.}\;
We show that $k < \sqrt{n}$:
\begin{alignat}{2}
    d &> \sqrt{n}
        &\qquad& \text{(assumption of this case)} \notag \\
    \frac{1}{d} &< \frac{1}{\sqrt{n}}
        &\qquad& \text{(taking reciprocals of positive quantities reverses the inequality)} \notag \\
    \frac{n}{d} &< \frac{n}{\sqrt{n}}
        &\qquad& \text{(multiplying both sides by $n > 0$)} \notag \\
    k = \frac{n}{d} &< \sqrt{n}
        &\qquad& \text{(simplifying $n/\sqrt{n} = \sqrt{n}$)} \notag
\end{alignat}

\medskip
\noindent In both cases at least one of $d$ or $k$ is less than or equal to
$\sqrt{n}$.
\end{proof}

\smallskip

\begin{remark}
This lemma tells us that every divisor $d$ of $n$ either is at most $\sqrt{n}$,
or its complementary divisor $n/d$ is at most $\sqrt{n}$.  In other words,
to enumerate \emph{all} divisors of $n$ it suffices to check only the
natural numbers up to $\sqrt{n}$.
\end{remark}

% ===============================================================
\section*{Main Theorem}
% ===============================================================

\begin{theorem}
For every $n \in \mathbb{N}$, the number of divisors $d(n)$ satisfies
\[
    d(n) \;\leq\; 2\sqrt{n}.
\]
\end{theorem}

\begin{proof}
Let $D$ denote the set of all positive divisors of $n$.  By definition
$|D| = d(n)$.

\smallskip
\noindent\textbf{Partition of $D$.}\;
We partition $D$ into three pairwise disjoint subsets:
\begin{align*}
    S_1 &= \bigl\{\, d \in D \;:\; d < \sqrt{n} \,\bigr\}, \\
    S_2 &= \bigl\{\, d \in D \;:\; d = \sqrt{n} \,\bigr\}, \\
    S_3 &= \bigl\{\, d \in D \;:\; d > \sqrt{n} \,\bigr\}.
\end{align*}
Since $S_1, S_2, S_3$ cover $D$ without overlap:
\[
    d(n) = |S_1| + |S_2| + |S_3|.
\]

% -----------------------------------------------------------
\bigskip
\noindent\textbf{Step 1: A bijection between $S_1$ and $S_3$.}

We claim $|S_1| = |S_3|$. We prove this by exhibiting a bijection
\[
    \varphi \;:\; S_1 \;\longrightarrow\; S_3, \qquad \varphi(d) = \frac{n}{d}.
\]
We verify in turn that $\varphi$ is well-defined, injective, and
surjective.\footnote{%
  Let $f\colon A\to B$ be a function.
  $f$ is \emph{injective} if $f(a)=f(b)\Rightarrow a=b$.\;
  $f$ is \emph{surjective} if $\forall\,b\in B\ \thickspace \exists\,a\in A$ such that
  $f(a)=b$.}

\bigskip
\noindent\textbf{\textit{Well-definedness.}}

\smallskip
\noindent
Let $d \in S_1$ be arbitrary; we must show $\varphi(d) = n/d \in S_3$.
Since $d \in S_1$, $d$ divides $n$, so $n/d \in \mathbb{N}$ and
$n/d$ divides $n$. It remains to check $n/d > \sqrt{n}$. Since $d < \sqrt{n}$,
\begin{alignat}{2}
    \frac{1}{d} &> \frac{1}{\sqrt{n}}
        &\qquad& \text{(taking reciprocals of positive quantities reverses the inequality)} \notag \\
    \frac{n}{d} &> \frac{n}{\sqrt{n}} = \sqrt{n}
        &\qquad& \text{(multiplying both sides by $n > 0$)} \notag
\end{alignat}
Hence $\varphi(d) \in S_3$ for all $d \in S_1$, and $\varphi$ is well-defined.

\bigskip
\noindent\textbf{\textit{Injectivity.}}

\smallskip
\noindent
Suppose $\varphi(a) = \varphi(b)$ for some $a, b \in S_1$, i.e.\ $n/a = n/b$.
\[
    \frac{n}{a} = \frac{n}{b}
    \;\implies\; n = \frac{na}{b}
    \;\implies\; bn = an
    \;\implies\; b = a.
\]
Hence $\varphi$ is injective.

\bigskip

\pagebreak
\noindent\textbf{\textit{Surjectivity.}}

\smallskip
\noindent\textit{Scratch work.}\;
\[
    \varphi(a) = b
    \;\implies\; \frac{n}{a} = b
    \;\implies\; n = ab
    \;\implies\; a = \frac{n}{b}.
\]

\noindent
Let $b \in S_3$ be arbitrary. Set $a := n/b$. Since $b \in S_3$,
$b$ divides $n$, so $a \in \mathbb{N}$ and $a$ divides $n$. Moreover,
since $b > \sqrt{n}$,
\[
    a = \frac{n}{b} < \frac{n}{\sqrt{n}} = \sqrt{n},
\]
so $a \in S_1$. Then
\[
    \varphi(a) = \frac{n}{a} = \frac{n}{\dfrac{n}{b}} = n \cdot \frac{b}{n} = b.
\]
Since $b \in S_3$ was arbitrary, $\varphi$ is surjective.

\bigskip
\noindent
Having verified that $\varphi$ is well-defined, injective, and surjective,
we conclude that $\varphi$ is a bijection from $S_1$ to $S_3$, and therefore
$|S_1| = |S_3|$.


% -----------------------------------------------------------
\bigskip
\noindent\textbf{Step 2: Bounding $|S_1|$.}

Using $|S_1| = |S_3|$, we rewrite the divisor count as
\[
    d(n) = |S_1| + |S_2| + |S_1|.
\]
Since $S_1$ consists of distinct naturals all strictly less than $\sqrt{n}$,
its elements are among $\{1, 2, \ldots, m\}$ where $m < \sqrt{n}$ is the
largest such natural (if $S_1 = \varnothing$, i.e.\ $n = 1$, then
$|S_1| = 0 < 1 = \sqrt{n}$ trivially); hence
\[
    |S_1| \leq m < \sqrt{n}.
\]

% -----------------------------------------------------------

\noindent\textbf{Step 3: Final bounds.}

We distinguish two cases.

\medskip
\noindent\textit{Case A: $n$ is not a perfect square.}

$\sqrt{n} \notin \mathbb{N}$, so no divisor of $n$ can equal $\sqrt{n}$; hence
$S_2 = \varnothing$ and $|S_2| = 0$.  Therefore
\[
    d(n) = |S_1| + 0 + |S_1| = 2|S_1| < 2\sqrt{n}.
\]

\medskip
\noindent\textit{Case B: $n$ is a perfect square.}

$\sqrt{n} \in \mathbb{N}$ and $\sqrt{n}$ divides $n$ (since
$\sqrt{n} \cdot \sqrt{n} = n$).  Thus $S_2$ contains exactly one element,
namely $\sqrt{n}$, so $|S_2| = 1$.  Therefore
\[
    d(n) = |S_1| + 1 + |S_1| = 2|S_1| + 1 < 2\sqrt{n} + 1.
\]
Since $n$ is a perfect square, $2\sqrt{n} \in \mathbb{N}$, so $d(n) < 2\sqrt{n}+1$
tightens to\footnote{For any two integers $a$ and $b$, $a < b + 1 \Leftrightarrow a \leq b$,
since there is no integer strictly between $b$ and $b+1$. Note that this would
fail for reals: $a = b + 0.5$ satisfies $a < b+1$ but not $a \leq b$.}
\[
    d(n) \leq 2\sqrt{n}.
\]

\medskip
\noindent In both cases we obtain $d(n) \leq 2\sqrt{n}$, which completes the
proof.
\end{proof}

\end{document}
